% !PRINTED 87
\weberpage{87}
Accordingly the congruence (17) has as a consequence:
\begin{gather}
e_{1,h}i_{1,h}p^{\pi-\pi_1}x_1 + e_{2,h}i_{2,h}p^{\pi-\pi_2}x_2 + \cdots e_{\varrho,h}i_{\varrho,h}p^{\pi-\pi_\varrho}x_\varrho \notag\\
\equiv 0\ (\mathrm{mod}\ p^\pi) \qquad h = 1, 2, \dots \mu, \tag{19}
\end{gather}
and the remaining primes dividing $J$ each yield just such a
system of congruences, which in their totality are completely
equivalent in meaning to the congruence (17). These
congruences (19) must not admit of being satisfied in any other way
than by the numbers
\[
p^{\pi-\pi_1}x_1,\ p^{\pi-\pi_2}x_2,\ \dots\ p^{\pi-\pi_\varrho}x_\varrho
\]
being divisible by $p^\pi$.

    If, then, we set for abbreviation:
\[
p^{\pi-\pi_1}x_1 = y_1,\ p^{\pi-\pi_2}x_2 = y_2,\ \dots\ p^{\pi-\pi_\varrho}x_\varrho = y_\varrho,
\]
then (19) takes the form:
\weq{20}{e_{1,h}i_{1,h}y_1 + e_{2,h}i_{2,h}y_2 + \cdots e_{\varrho,h}i_{\varrho,h}y_\varrho \equiv 0\ (\mathrm{mod}\ p^\pi),}
and these congruences are to have
\[
y_1 \equiv 0,\ y_2 \equiv 0,\ \dots\ y_\varrho \equiv 0\ (\mathrm{mod}\ p^\pi)
\]
as a necessary consequence.

    The system of linear congruences (20) can now be discussed
by the same means as a system of linear equations, namely by
means of determinants (cf. Chapter II of the first volume).

    For a better overview we shall write the congruences (20)
for the moment thus:
\begin{gather}
a_{1,1}y_1 + a_{2,1}y_2 + \cdots + a_{\varrho,1}y_\varrho \equiv 0 \notag\\
a_{1,2}y_1 + a_{2,2}y_2 + \cdots + a_{\varrho,2}y_\varrho \equiv 0\ (\mathrm{mod}\ p^\pi). \tag{21}\\
\dots\dots\dots\dots\dots\dots\dots\dots\dots\dots \notag\\
a_{1,\mu}y_1 + a_{2,\mu}y_2 + \cdots + a_{\varrho,\mu}y_\varrho \equiv 0 \notag
\end{gather}
If the coefficients $a_{k,h}$ are all divisible by $p$, then
these congruences are already satisfied when the $y$ are all only
divisible by $p^{\pi-1}$, and this case must therefore not occur.

    We shall now assume that from the matrix
\weq{22}{
\begin{matrix}
a_{1,1}, & a_{2,1}, & \dots & a_{\varrho,1} \\
a_{1,2}, & a_{2,2}, & \dots & a_{\varrho,2} \\
\dots & \dots & \dots & \dots \\
a_{1,\mu}, & a_{2,\mu}, & \dots & a_{\varrho,\mu}
\end{matrix}
}
%CONT>>
